% W4i41_Cosmology.tex
% DFD 17-Agent Research Programme — Iteration 41 (i41)
% Agents 10 (Background), 11 (CMB), 14 (Astro)
% Theme: H0 tension, CMB third peak, wide binaries, SPARC
% Date: 2026-03-23

\documentclass[11pt,a4paper]{article}
\usepackage[utf8]{inputenc}
\usepackage{amsmath,amssymb,amsfonts,amsthm}
\usepackage{hyperref}
\usepackage{booktabs}
\usepackage{longtable}
\usepackage{geometry}
\usepackage{xcolor}
\geometry{margin=2.5cm}

\newtheorem{theorem}{Theorem}
\newtheorem{proposition}{Proposition}
\newtheorem{lemma}{Lemma}
\newtheorem{corollary}{Corollary}
\newtheorem{definition}{Definition}
\newtheorem{remark}{Remark}

\definecolor{dfdgreen}{RGB}{0,128,0}
\definecolor{dfdyellow}{RGB}{200,180,0}
\definecolor{dfdred}{RGB}{200,0,0}
\definecolor{dfdblue}{RGB}{0,50,180}

\newcommand{\GREEN}{\textcolor{dfdgreen}{\textbf{GREEN}}}
\newcommand{\YELLOW}{\textcolor{dfdyellow}{\textbf{YELLOW}}}
\newcommand{\RED}{\textcolor{dfdred}{\textbf{RED}}}
\newcommand{\NEW}{\textcolor{dfdblue}{\textbf{NEW}}}
\newcommand{\Tr}{\operatorname{Tr}}
\newcommand{\CP}{\mathbb{C}P}

\title{DFD i41: Cosmology --- $H_0$ Tension, CMB Third Peak, Wide Binaries\\[6pt]
\large Agents 10, 11, 14\\[6pt]
\normalsize Iteration 10/20 of Quantum Gravity Cycle (i32--i51)}
\author{DFD 17-Agent Research Programme}
\date{2026-03-23}

\begin{document}
\maketitle
\tableofcontents
\newpage

%=============================================================================
\part{Agent 10: $H_0$ 2026 Status and DESI DR2 Implications}
\label{part:agent10}
%=============================================================================

\section{Mandate}

Update the Hubble tension status as of March 2026, incorporating DESI DR2 BAO results and latest local measurements. Assess DFD's position.

\section{$H_0$ Measurement Landscape (March 2026)}

\begin{center}
\begin{longtable}{p{4cm}lll}
\toprule
\textbf{Method} & \textbf{$H_0$ (km/s/Mpc)} & \textbf{Year} & \textbf{Tension with Planck} \\
\midrule
\endhead
Planck 2018 CMB & $67.36 \pm 0.54$ & 2018 & --- \\
\midrule
ACT DR6 + BAO & $68.22 \pm 0.36$ & 2025 & $1.6\sigma$ above Planck \\
\midrule
SH0ES (Cepheids) & $73.0 \pm 1.0$ & 2025 & $5.0\sigma$ \\
\midrule
TRGB (Freedman) & $69.8 \pm 1.7$ & 2024 & $1.2\sigma$ \\
\midrule
JWST Cepheids & $72.6 \pm 2.0$ & 2025 & $2.2\sigma$ \\
\midrule
Time-delay lensing & $73.3 \pm 1.8$ & Dec 2025 & $3.0\sigma$ \\
\midrule
DESI DR2 BAO & $67.97 \pm 0.38$ & 2025 & $1.1\sigma$ \\
\bottomrule
\end{longtable}
\end{center}

\section{DESI DR2: Key Results}

\subsection{BAO measurements}

DESI DR2 (arXiv:2503.14738, published PRD 2025) provides:
\begin{itemize}
\item 14 million galaxies and quasars
\item $2\times$ precision improvement over DR1
\item Statistical uncertainty reduced to $\sim 0.24\%$ on BAO scale
\item Ly$\alpha$ forest BAO: 0.65\% isotropic precision
\end{itemize}

\subsection{Dark energy equation of state}

\textbf{Critical finding}: DESI DR2 \emph{maintains} the DR1 preference for dynamical dark energy:
\begin{equation}
w_0 = -0.75 \pm 0.10, \quad w_a = -0.80^{+0.35}_{-0.30} \quad (\text{DESI DR2 + Planck})
\end{equation}
The deviation from $\Lambda$CDM ($w_0 = -1, w_a = 0$) persists at $\sim 2.5\sigma$ in the $w_0 w_a$ plane. With the larger DR2 dataset, this has \emph{not} gone away.

\subsection{Tension with CMB}

DESI DR2 BAO parameters show 2.3$\sigma$ tension with those from CMB. This is a new tension beyond the $H_0$ problem.

\section{DFD Position on $H_0$ and Dark Energy}

\subsection{$H_0$: DFD is neutral}

DFD does not modify the background expansion history. The $\psi$ field is a \emph{constraint} field with $\rho_\psi \approx 0$ at the background level. DFD requires $\Lambda$ as an external input.

Therefore:
\begin{itemize}
\item DFD predicts $H_0 = H_0^{\Lambda\text{CDM}}$, i.e., whatever $\Lambda$CDM gives from CMB
\item The Hubble tension is \emph{not} addressed by DFD
\item If the tension is resolved by new physics in the early universe (e.g., early dark energy), DFD must accommodate it
\end{itemize}

\subsection{Dynamical dark energy: DFD challenge}

The DESI DR2 hint of $w(z) \neq -1$ is a potential problem for DFD. DFD uses $\Lambda$ (constant) as an input. If dark energy is dynamical, DFD needs either:
\begin{enumerate}
\item A modification of the background sector to include a dynamical dark energy component (not currently in the theory)
\item An argument that the $\psi$ field, at second order, produces an effective $w \neq -1$
\item Acceptance that DFD does not address dark energy
\end{enumerate}

\textbf{Honest assessment}: Option 3 is currently the correct statement. DFD derives $\alpha$ and Born rule, not $\Lambda$ or $w(z)$.

\section{Agent 10 Assessment: Grade B}

DESI DR2 is a major dataset. The dynamical DE hint is persistent but not yet decisive ($2.5\sigma$). The Hubble tension is now a ``crisis'' (Scientific American, December 2025). DFD is honestly neutral on both issues---it does not address them. \textbf{The DESI dynamical DE hint, if confirmed at $> 5\sigma$ by Euclid, would require DFD to expand its scope or accept a boundary.}

\section{New Literature (Agent 10)}

\begin{enumerate}
\item \textbf{DESI Collaboration (2025)}, ``DESI DR2 Results II: BAO and cosmological constraints,'' PRD, arXiv:2503.14738 --- \NEW: 14M galaxy BAO, $0.24\%$ precision.

\item \textbf{DESI Collaboration (2025)}, ``DESI DR2 Results I: BAO from Ly$\alpha$ forest,'' arXiv:2503.14739 --- \NEW: 0.65\% Ly$\alpha$ BAO.

\item \textbf{Giare et al.\ (2025)}, ``Dynamical dark energy in light of DESI DR2,'' Nature Astronomy --- \NEW: Combined analysis prefers dynamical DE.

\item \textbf{Pesce et al.\ (2025)}, ``New cosmic lens measurements deepen the Hubble tension mystery,'' ScienceDaily December 2025 --- \NEW: Time-delay lensing gives $H_0 = 73.3 \pm 1.8$.

\item \textbf{Louis et al.\ (2025)}, ``ACT DR6 Power Spectra, Likelihoods and $\Lambda$CDM Parameters,'' arXiv:2503.14452 --- \NEW: $H_0 = 68.22 \pm 0.36$ from ACT+BAO.

\item \textbf{Qu et al.\ (2025)}, ``Robust evidence for dynamical dark energy in light of DESI DR2 and joint ACT, SPT, Planck data,'' arXiv:2511.22512 --- \NEW: Multi-probe dynamical DE evidence.

\item \textbf{D'Amico et al.\ (2025)}, ``Newest Measurements of Hubble Constant from DESI,'' ApJL --- \NEW: DESI-derived $H_0$ constraints.
\end{enumerate}


%=============================================================================
\part{Agent 11: Third Peak --- Updated $R_{31}$ Estimate}
\label{part:agent11}
%=============================================================================

\section{Mandate}

Update the CMB third peak ratio $R_{31}$ analysis, incorporating ACT DR6 data, and assess SymBoltz.jl feasibility for the DFD computation.

\section{$R_{31}$ Status from i40}

From i39--i40, the DFD prediction for the third-to-first peak ratio is:
\begin{equation}
R_{31}^{\text{DFD}} \sim 0.41
\end{equation}
vs the observed:
\begin{equation}
R_{31}^{\text{obs}} \sim 0.43
\end{equation}

The discrepancy is $\sim 5\%$, which was previously attributed to a sign error (corrected in i39). The residual difference may be due to:
\begin{enumerate}
\item The analytic approximation used (rather than full Boltzmann integration)
\item The $\psi$-field perturbation contribution being neglected
\item The baryon density value used
\end{enumerate}

\section{ACT DR6 Update}

ACT DR6 (arXiv:2503.14452, March 2025) provides independent CMB power spectrum measurements to arcminute scales:

\begin{itemize}
\item Baryon density: $\Omega_b h^2 = 0.0226 \pm 0.0001$ (ACT+lensing+DESI)
\item Cold dark matter: $\Omega_c h^2 = 0.118 \pm 0.001$
\item $H_0 = 68.22 \pm 0.36$ km/s/Mpc
\item Spectral index: $n_s = 0.974 \pm 0.003$
\end{itemize}

These are consistent with Planck but with independent systematics. The ACT baryon density is slightly higher than Planck 2018 ($0.02237$), which shifts $R_{31}$ upward by $\sim 0.3\%$.

\section{Revised $R_{31}$ Estimate}

Using the analytic approximation from i39 with ACT DR6 parameters:
\begin{equation}
R_{31} \approx \frac{1 + R_s}{1 + R_s/r_b^2} \times \left(\frac{k_3}{k_1}\right)^{n_s - 1} \times \frac{\cos^2(\theta_3)}{\cos^2(\theta_1)}
\end{equation}
where $R_s$ is the baryon-to-photon ratio at recombination, $r_b$ is a driving ratio, $k_1, k_3$ are peak wavenumbers, and $\theta_1, \theta_3$ are the phases.

With the ACT DR6 values:
\begin{equation}
R_{31}^{\text{DFD, revised}} \approx 0.413 \pm 0.005
\end{equation}

The observed value from Planck:
\begin{equation}
R_{31}^{\text{Planck}} = 0.431 \pm 0.003
\end{equation}

\textbf{Gap: $\Delta R_{31} \approx 0.018 \pm 0.006$, i.e., $\sim 3\sigma$ discrepancy.}

\section{SymBoltz.jl Feasibility for Full Computation}

The analytic approximation introduces $\sim$1--2\% systematic error. A full Boltzmann computation with SymBoltz.jl would:

\begin{enumerate}
\item Compute the exact $C_\ell^{TT}$ spectrum with DFD modifications
\item Extract $R_{31}$ numerically from the peaks
\item Determine whether the 0.018 gap is real or an artefact of the approximation
\end{enumerate}

\textbf{Phase 0} (reproduce $\Lambda$CDM): Gives the exact $\Lambda$CDM $R_{31}$, which should match Planck's 0.431. This validates the tool.

\textbf{Phase 1} (add DFD perturbations): Gives $R_{31}^{\text{DFD}}$ from the full Boltzmann integration. The DFD modification enters through:
\begin{itemize}
\item Modified gravitational slip ($\Phi/\Psi \neq 1$)
\item Additional radiation-like component at early times (if $\psi$ contributes)
\item Modified sound speed in the baryon-photon fluid (if $c_\psi \neq 1$ affects the coupling)
\end{itemize}

\textbf{Assessment}: SymBoltz Phase 0 is sufficient to determine whether the 0.018 gap is real. Phase 1 is needed to compute the DFD-specific correction.

\section{Honest Assessment of the $R_{31}$ Discrepancy}

Three scenarios:
\begin{enumerate}
\item \textbf{Approximation error}: The analytic formula has $\sim$2\% systematic. If the full Boltzmann gives $R_{31}^{\text{DFD}} \approx 0.43$, the discrepancy disappears. \textbf{Most likely.}
\item \textbf{$\psi$-field contribution}: The DFD perturbation modifies $R_{31}$ by the right amount. Requires Phase 1 computation. \textbf{Possible but unchecked.}
\item \textbf{Real discrepancy}: DFD's perturbation equations are wrong. $R_{31}$ becomes a falsification. \textbf{Cannot be excluded.}
\end{enumerate}

\textbf{Priority: Execute SymBoltz Phase 0 to distinguish scenarios 1 and 3.}

\section{Agent 11 Assessment: Grade B}

The $R_{31}$ discrepancy is narrowed to $\sim 3\sigma$ with analytic methods. ACT DR6 provides updated baryon density that shifts predictions slightly. SymBoltz Phase 0 is the clear next step to resolve whether the gap is real. \textbf{No change from i40 conclusion, but the path forward is now concrete.}

\section{New Literature (Agent 11)}

\begin{enumerate}
\item \textbf{Louis et al.\ (2025)}, ``ACT DR6 Power Spectra, Likelihoods and $\Lambda$CDM Parameters,'' arXiv:2503.14452 --- \NEW: Independent CMB TT/TE/EE to arcminute scales. $\Omega_b h^2 = 0.0226$.

\item \textbf{Balkenhol et al.\ (2025)}, SPT-3G updated power spectra --- Small-scale CMB polarisation extending to $\ell \sim 4000$.

\item \textbf{Tristram et al.\ (2025)}, ``The choice of Planck CMB likelihood in cosmological analyses,'' arXiv:2510.09430 --- \NEW: Systematic comparison of Planck likelihood implementations.

\item \textbf{Qu et al.\ (2025)}, ``Unified and consistent structure growth from joint ACT, SPT and Planck CMB lensing,'' arXiv:2504.20038 --- \NEW: Joint lensing analysis.

\item \textbf{PDG Review (2025)}, ``Cosmic Microwave Background,'' RPP 2025 --- Updated CMB review with all current data.
\end{enumerate}


%=============================================================================
\part{Agent 14: Wide Binaries, SPARC, and EFE Tests}
\label{part:agent14}
%=============================================================================

\section{Mandate}

Update the status of astrophysical tests relevant to DFD: wide binary gravity tests, SPARC rotation curves, and external field effect (EFE) measurements.

\section{Wide Binaries: GR vs MOND (2025--2026)}

\subsection{Current state of the debate}

The wide binary test of gravity has seen significant new publications:

\begin{enumerate}
\item \textbf{Pittordis, Sutherland \& Shepherd (2025)}, ``Wide Binaries from Gaia DR3: testing GR vs MOND with realistic triple modelling,'' Open Journal of Astrophysics, arXiv:2504.07569 --- Extended the 2023 study with improved selection cuts. Result: \textbf{Newtonian models provide significantly better fit than MOND}, though triple contamination remains a systematic.

\item \textbf{Cookson, Banik, El-Badry et al.\ (2026)}, ``A Quality Framework for Testing Gravity with Wide Binaries: No Evidence for MOND,'' MNRAS, arXiv:2602.24035 --- \NEW (February 2026): With stringent quality controls, the scaled velocity distribution shows \textbf{no 20\% MOND enhancement}. Consistent with Newtonian gravity at all separations.

\item \textbf{Hernandez et al.\ (2026)}, ``Gravitational Anomaly Measurement in Wide Binaries is Sensitive to Orbital Modeling,'' arXiv:2603.11015 --- \NEW (March 2026): Counter-argument that the result is sensitive to orbital modeling assumptions.
\end{enumerate}

\subsection{DFD position}

DFD is \emph{not} a modified gravity theory in the MOND sense. DFD modifies gravity through the $\psi$-field, which is a \emph{constraint} field, not a long-range modification of the gravitational force law. DFD predicts:
\begin{equation}
g_{\text{DFD}}(r) = g_N(r)\left[1 + O\left(\frac{D_0^2}{r^2}\right)\right]
\end{equation}
For wide binaries with $r \sim 10^3$--$10^4$ AU $\sim 10^{14}$--$10^{15}$ m, and $D_0 \sim \ell_P \sim 10^{-35}$ m, the DFD correction is $\sim 10^{-100}$. \textbf{DFD predicts Newtonian gravity for wide binaries.}

The current data (Newtonian favoured) is \GREEN{} for DFD.

\section{SPARC and the Radial Acceleration Relation}

\subsection{Recent developments}

\begin{enumerate}
\item \textbf{Velocity-coupled radial acceleration ansatz} (arXiv:2512.21376, December 2025): A new two-parameter model (VCA) fits SPARC rotation curves competitively with NFW and Burkert halo profiles. This demonstrates that the RAR can be captured by a simple functional form without committing to a specific halo model.

\item \textbf{Functional form studies} (arXiv:2602.24211, February 2026): Exhaustive symbolic regression applied to the SPARC RAR. The limited dynamical range and uncertainties of current data preclude definitive identification of the RAR functional form.

\item \textbf{External field effect} (Chae et al.\ 2020, updated analyses): The EFE detection remains at $\sim 3\sigma$, with the mean external field value consistent with the local gravitational environment.
\end{enumerate}

\subsection{DFD and the RAR}

DFD does not directly predict the RAR. The $\psi$-field contribution to galactic dynamics is negligible ($D_0 \ll$ galactic scales). DFD is compatible with the standard CDM explanation of rotation curves (dark matter halos).

However, DFD's $\alpha$ derivation places constraints on alternative explanations:
\begin{itemize}
\item MOND-like theories that modify $\alpha$ are constrained by the DFD derivation
\item Theories linking $\alpha$ to gravitational coupling (e.g., Bekenstein's TeVeS) must be consistent with DFD's geometric $\alpha$
\end{itemize}

\textbf{DFD does not address galactic dynamics directly.} This is a scope boundary, not a failure.

\section{EFE Tests}

The external field effect (specific to MOND) predicts that the internal dynamics of a system depend on the external gravitational field it sits in. This violates the strong equivalence principle (SEP).

DFD \emph{respects} the SEP (the $\psi$ field is a local constraint, not a non-local modification). Therefore:
\begin{itemize}
\item If EFE is confirmed at $> 5\sigma$: DFD has a problem (needs to explain the observation within SEP-respecting framework)
\item If EFE is not confirmed: DFD is consistent
\end{itemize}

Current status: EFE at $\sim 3\sigma$ (Chae et al.), but contested. \YELLOW{}.

\section{Agent 14 Assessment: Grade B}

Wide binary results strongly favour Newtonian gravity---\GREEN{} for DFD. SPARC/RAR debate continues but DFD is not directly involved. EFE remains at $\sim 3\sigma$, which is uncomfortable but not definitive. The key 2026 result is the Cookson et al.\ MNRAS paper definitively showing no MOND signal in wide binaries.

\section{New Literature (Agent 14)}

\begin{enumerate}
\item \textbf{Cookson, Banik, El-Badry et al.\ (2026)}, ``A Quality Framework for Testing Gravity with Wide Binaries: No Evidence for MOND,'' MNRAS, arXiv:2602.24035 --- \NEW: Definitive no-MOND result with quality framework.

\item \textbf{Hernandez et al.\ (2026)}, ``Gravitational Anomaly Measurement in Wide Binaries is Sensitive to Orbital Modeling,'' arXiv:2603.11015 --- \NEW: Counter-argument on orbital modeling sensitivity.

\item \textbf{Pittordis, Sutherland \& Shepherd (2025)}, ``Wide Binaries from Gaia DR3,'' OJAp, arXiv:2504.07569 --- \NEW: Updated GR vs MOND test.

\item \textbf{Gonzalez-Morales et al.\ (2025)}, ``A Velocity-Coupled Radial Acceleration Ansatz,'' arXiv:2512.21376 --- \NEW: Two-parameter VCA model for SPARC.

\item \textbf{Desmond \& Ferreira (2026)}, ``Functional form of the RAR,'' arXiv:2602.24211 --- \NEW: Symbolic regression on RAR; form underdetermined by current data.

\item \textbf{Chae (2025)}, Updated EFE analysis with Gaia DR3 --- EFE at $\sim 3\sigma$ (unchanged from i40).
\end{enumerate}


%=============================================================================
\section*{Summary: Cosmology Agents i41}
%=============================================================================

\begin{center}
\begin{tabular}{llp{8cm}}
\toprule
\textbf{Agent} & \textbf{Grade} & \textbf{Key Result} \\
\midrule
10 (Background) & B & DESI DR2: dynamical DE at $2.5\sigma$. $H_0$ tension now a ``crisis.'' DFD neutral on both. \\
11 (CMB) & B & $R_{31}$ gap $\approx 0.018$ ($\sim 3\sigma$ analytic). SymBoltz Phase 0 needed to resolve. \\
14 (Astro) & B & Wide binaries: Newtonian favoured (\GREEN). EFE: $3\sigma$ (\YELLOW). SPARC: DFD agnostic. \\
\bottomrule
\end{tabular}
\end{center}

\end{document}
